\documentclass[11pt,a4paper]{article}
\usepackage[utf8]{inputenc}
\usepackage[T1]{fontenc}
\usepackage[margin=1.35cm]{geometry}
\usepackage{enumitem}
\usepackage{xcolor}
\usepackage{parskip}
\definecolor{gris}{HTML}{1F2937}
\begin{document}
\pagestyle{empty}
\begin{center}
{\Large\textbf{Gustavo Vega}}\\
Operaciones de tienda y coordinación de equipos\\
gusvegabol@gmail.com \quad 669 549 933 \quad linkedin.com/in/gusvegabol
\end{center}
\section*{Perfil profesional}
Profesional de operaciones con experiencia directa en supermercados, coordinación de equipos y mejora de procesos. En Herfrailes S. L. organicé el trabajo de tres tiendas, apoyé la gestión de responsables y conecté ventas, pedidos, stock, logística y atención al cliente. Busco aportar autonomía, responsabilidad y orientación a resultados como Adjunto/a de tienda.
\section*{Propuesta de valor}
Combino visión operativa, trato directo con clientes y capacidad para convertir información en tareas claras. He diseñado sistemas para prever pedidos, redistribuir stock y seguir el trabajo diario, manteniendo el foco en el servicio, el equipo y la rentabilidad.
\section*{Experiencia profesional relevante}
\textbf{2011--2024 | Director Ejecutivo -- Herfrailes S. L.} Coordiné responsables de tienda y de sección, organicé equipos polivalentes entre frutas, charcutería, caja y reposición, atendí reclamaciones de clientes e implanté Trello para asignar y seguir tareas operativas.

\textbf{2011--2024 | Operaciones y logística -- Herfrailes S. L.} Diseñé una base de datos de ventas y un algoritmo de pedido que redujo un 30\% las caducidades y un 80\% los productos sin venta durante más de un mes. Organicé la redistribución de stock entre tres tiendas y el seguimiento del mantenimiento de vehículos.

\textbf{2011--2024 | Gestión de personas y comunicación -- Herfrailes S. L.} Estructuré tareas, documentación y canales de comunicación para responsables y equipos. El seguimiento mediante Trello llegó a gestionar hasta 200 tareas pendientes de revisión diaria.

\textbf{2009--2011 | Programador -- INERZA, S. A.} Desarrollé soluciones de gestión de vivienda pública y generación de documentación.

\textbf{1985--2009 | Analista programador y vendedor -- Computerland, General de Software de Canarias y otras etapas} Desarrollé soluciones de gestión y contabilidad, vendí hardware y servicios, impartí formación y mantuve trato directo con clientes.
\section*{Competencias y herramientas}
\begin{itemize}[leftmargin=*]
\item Operaciones de tienda, organización de equipos y seguimiento de tareas
\item Ventas, previsión de pedidos, stock y logística entre tiendas
\item Orientación al cliente, resolución de reclamaciones y comunicación
\item Excel y Word avanzados; automatización y mejora de procesos
\end{itemize}
\section*{Formación}
Bachillerato -- título oficial expedido por el Ministerio de Educación y Ciencia\\
Formación Profesional, Técnico Administrativo -- IFP Santa María de Guía\\
Máster en Inteligencia Artificial con Certificación Universitaria -- BIG school / Universidad Isabel I, en curso
\section*{Información adicional}
Experiencia específica en supermercados, logística local y atención de reclamaciones. Disponibilidad para turnos: pendiente de confirmar.
\end{document}
